\subsection{Puntos únicos de falla y riesgos residuales}
\label{sec:spof}

RT-02.11 exige identificar las dependencias singulares que subsisten y justificar el riesgo residual. La Tabla~\ref{tab:spof} distingue su efecto, la continuidad prevista y la condición de aceptación. Una réplica o una cola reduce el impacto, pero no elimina por declaración la dependencia.

\begin{tablalafrox}{Dependencias singulares y riesgo residual}{tab:spof}{F{0.21} F{0.26} F{0.27} Y}{\cab{Dependencia} & \cab{Efecto si falla} & \cab{Continuidad prevista} & \cab{Aceptación}}
Identidad central & No hay altas ni revocación remota. & Asignaciones vigentes y servicio local. & Prueba de relevo 24 h. \\
ERP y emisión DTE & No se confirma nueva guía. & Se conserva trabajo preparado; salida afectada se bloquea. & Contingencia tributaria validada. \\
Maestro bodega Talca & Se pierde consolidación y reserva global. & Cada sitio opera localmente y reconcilia. & Conflicto de stock declarado. \\
SII y pasarela de pago & No hay respuesta en línea. & Reintento y cobro alternativo aprobado. & Pendiente no equivale a aprobado. \\
Planificador experto & Faltan decisiones ante excepciones. & Reglas de M4 y suplente capacitado. & Prueba antes de su retiro. \\
\end{tablalafrox}

{\footnotesize Fuente: elaboración propia de LafroX a partir de Bases Técnicas Transversales, RT-02.11, y caso 02.\par}

Las dependencias externas conservan riesgo residual que exige prueba o contingencia aprobada.

La dependencia del ERP y del SII no es aceptable para un nuevo despacho sin el mecanismo tributario de contingencia acordado. El riesgo de identidad local tampoco se declara cerrado antes de probar cada relevo de turno durante un corte de 24 horas. El emplazamiento y la redundancia física se verifican en 4.2.

